\chapter{Présentation de la structure d'accueil}

\begin{simplebox2}
Dans cette partie, nous allons présenter brièvement la structure ayant servi de cadre de stage c’est-à-dire l’Agence Nationale de la Statistique et de la Démographie (ANSD). Pour ce faire, nous ferons ressortir ses missions et attributs après quoi nous présenterons les différentes directions qui la composent et mettrons le focus sur la direction de la Méthodologie, de la Coopération statistique et de l'Innovation.
\end{simplebox2}

\section{Historique, rôle et mission de l'ANSD}

Créée par la loi n°2004-21 du 21 juillet 2004, portant organisation des activités statistiques au sein du Système Statistique National (SSN), l’Agence nationale de la Statistique et de la Démographie (ANSD) est une structure administrative dotée de la personnalité juridique et d'une autonomie de gestion. Elle est placée sous la tutelle du ministre de l'économie du plan et de la coopération. L’organisation administrative de l'ANSD a été définie par le décret n°2005-436 du 23 mai 2005.\\
l
La mission de l’ANSD s’articule autour des points suivants :

\begin{simplebox}
\begin{itemize}
    \item veiller à l’élaboration et à la mise en œuvre des programmes pluriannuels et annuels de statistique ;
    \item assurer la mise en application des méthodes, concepts, définitions, normes, classifications et nomenclatures approuvés par le Comité technique des programmes statistiques ;
    \item préparer les dossiers à soumettre aux réunions du Conseil national de la statistique et du Comité technique des programmes statistiques ;
    \item assurer le secrétariat et l’organisation des réunions du Conseil national de la statistique et du Comité technique des programmes statistiques ainsi que de ses sous-comités sectoriels ;
    \item réaliser des enquêtes d’inventaire à couverture nationale notamment les recensements généraux de la population et ceux des entreprises ;
    \item produire les comptes nationaux ;
    \item suivre la conjoncture et la prévision économiques en rapport avec le service en charge de la prévision et de la conjoncture économique ;
    \item élaborer et gérer les fichiers des entreprises et des localités ;
    \item élaborer les indicateurs économiques, sociaux et démographiques ;
    \item centraliser et diffuser les synthèses des données statistiques produites par l’ensemble du système statistique national ;
    \item favoriser le développement des sciences statistiques et la recherche économique appliquée relevant de sa compétence ;
    \item promouvoir la formation du personnel spécialisé pour le fonctionnement du système national d’information statistique par l’organisation des cycles de formation appropriés, notamment au sein d’une école à vocation régionale ou sous-régionale intégrée à l’agence.
\end{itemize}
\end{simplebox}

Autrement dit, l'ANSD a pour mission d'assurer la cohésion technique des activités du SSN. Elle coordonne les activités statistiques nationales et se charge de la production et de la diffusion des données statistiques. Ces données sont mises à disposition pour répondre aux besoins du gouvernement, des administrations publiques, du secteur privé, des partenaires au développement et du grand public. 

En outre, l’Agence est chargée du suivi de la coopération technique internationale en matière statistique. A ce titre, elle représente le Sénégal dans les instances sous-régionales, régionales et internationales relatives aux questions relevant de sa compétence et suit les activités des organisations internationales en ce qui concerne les questions statistiques.





%====================================================

\section{Structure organisationnelle}

La Direction générale comprend : le Directeur général, le Directeur général adjoint, les Conseillers techniques, l’Auditeur interne, la Cellule du contrôle de gestion, la Cellule de passation des marchés, la Cellule de communication, la Cellule juridique, la Cellule de gestion des partenariats, le Service de sécurité, le Service médical 
ainsi que les secrétariats du Directeur général et du Directeur général adjoint. 
Un Bureau de l’archivage et du courrier y est rattaché pour la gestion des archives et des courriers à l’arrivée et au départ de l’Agence.\\\\
En dehors de la Direction générale, l’ANSD compte huit directions à savoir :

\begin{simplebox}
\begin{itemize}
    \item La Direction des Statistiques Économiques et de la Comptabilité nationale (DSECN) ;
    \item La Direction des Statistiques Démographiques et Sociales (DSDS) ;
    \item La Direction des Systèmes d’Information et de la Diffusion (DSID) ;
    \item La Direction de l’Administration Générale et des Ressources Humaines (DAGRH) ;
    \item La Direction de l’École Nationale de la Statistique et de l’Analyse Économique (ENSAE) ;
    \item La Direction de la Méthodologie, de la Coopération statistique et de l’Innovation (DMCI) ;
    \item La Direction de l’Action Régionale (DAR) ;
    \item L’Agence comptable.
\end{itemize}
\end{simplebox}




%====================================================
\section{Présentation de la DMCI}

La \textbf{Direction de la Méthodologie, de la Coordination Statistique et de l’Innovation (DMCI)} 
est une direction récente. Elle a remplacé la Direction de la Méthodologie, de la Coordination 
statistique et des Partenariats (DMCP). Ce changement de dénomination est survenu à l’issue 
d’une réorganisation de l’organigramme de l’Agence qui a eu lieu le 15 juin 2023. 

Ces révisions se justifient par la perpétuelle quête de l’ANSD à s’arrimer aux standards 
internationaux et à s’orienter vers l’innovation. La DMCI élabore et vulgarise les méthodologies 
et les bonnes pratiques en matière statistique. Elle est notamment chargée d’harmoniser les 
concepts, nomenclatures et classifications utilisés à l’ANSD et au sein de tout le Système 
Statistique National (SSN). 

En outre, elle est responsable de la mise en place, progressive et coopérative, d’un cadre 
national d’assurance qualité, d’assurer la veille technologique ainsi que de promouvoir la 
recherche et l’innovation. La DMCI assure également la coordination et la programmation 
des activités statistiques du SSN et supervise la coopération internationale de l’ANSD. 
Elle sert d’interlocuteur entre l’ANSD et les autres structures membres du SSN. 

Elle appuie le Directeur général dans la préparation des dossiers relatifs au Conseil national 
de la Statistique (CNS) et au Comité technique des Programmes statistiques (CTPS), notamment 
la préparation des sessions, commissions et groupes de travail, ainsi que le suivi de l’exécution 
de leurs décisions. Elle veille au maintien d’un niveau élevé de maîtrise des techniques 
statistiques et à la gestion de la qualité des résultats produits. 


\begin{simplebox1}
La DMCI est structurée en deux divisions principales :  
\begin{itemize}
    \item La Division de la Méthodologie et de l’Innovation (DMI) ;
    \item La Division de la Coopération et de la Coordination Statistique (DCC).
\end{itemize}
\end{simplebox1}

\subsection{Division de la Méthodologie et de l’Innovation (DMI)}

La DMI est en charge de l’élaboration des méthodologies, nomenclatures et des classifications. 
Elle assure la mise en place d’un système d’assurance qualité et promeut la veille technologique 
ainsi que la recherche et l’innovation. Elle est composée de deux bureaux :


\begin{simplebox}
\begin{itemize}
    \item \textbf{Le Bureau des Méthodes et Normes (BMN)} : responsable de l’élaboration et 
    de la vulgarisation des méthodologies et normes, de la mise en œuvre du Cadre National 
    d’Assurance Qualité (CNAQ) et de la coordination de l’harmonisation méthodologique et 
    de la gestion des visas statistiques.
    \item \textbf{Le Bureau de la Recherche et de l’Innovation (BRI)} : en charge de la veille 
    technologique, de la promotion de l’innovation ainsi que de la coordination d’études 
    approfondies et de la production de rapports.
\end{itemize}
\end{simplebox}

\subsection{Division de la Coopération et de la Coordination Statistique (DCC)}

    
La DCC élabore les Stratégies Nationales de Développement de la Statistique (SNDS) et 
les Plans Stratégiques de Développement (PSD). Elle administre la plateforme PTA, réalise 
la cartographie du SSN et gère la coopération internationale. Elle comprend deux bureaux :

\begin{simplebox}
\begin{itemize}
    \item \textbf{Le Bureau de la Coordination Statistique (BCS)} dont les principales missions sont :
    \begin{itemize}
        \item Élaborer et assurer le suivi et l’évaluation des SNDS et des PSD ;
        \item Administrer la plateforme PTA et réaliser périodiquement la cartographie du SSN ;
        \item Animer le Système Statistique National avec l’organisation annuelle du CNS, 
        des réunions du CTPS et des sous-comités ;
        \item Administrer et animer le portail web du SSN et réaliser chaque année 
        l’enquête de satisfaction des utilisateurs ;
        \item Identifier et mettre en œuvre, en collaboration avec l’ENSAE ou d’autres structures, 
        des renforcements de capacités au bénéfice du SSN.
    \end{itemize}

    \item \textbf{Le Bureau de la Coopération Internationale (BCI)} dont les tâches consistent à :
    \begin{itemize}
        \item Suivre les indicateurs ODD et l’adhésion aux normes internationales (NSDD, NSDD Plus) ;
        \item Gérer les partenariats internationaux avec des institutions comme UNSD, PARIS21, et AFRISTAT.
    \end{itemize}
\end{itemize}
\end{simplebox}



%====================================================
\section{Bilan du Stage}
A  remplir
